% ============================================================================
% LANGENTWURF LE-009: Datenschutz & Privatsphäre
% Profil: S (Senioren 65+) | Tier: 2 (Standard)
% ============================================================================

\documentclass[11pt,a4paper]{article}
\usepackage[utf8]{inputenc}
\usepackage[T1]{fontenc}
\usepackage[ngerman]{babel}
\usepackage{geometry}
\usepackage{fancyhdr}
\usepackage{lastpage}
\usepackage{xcolor}
\usepackage{tcolorbox}
\usepackage{enumitem}
\usepackage{longtable}
\usepackage{hyperref}
\usepackage{amssymb}

\geometry{left=2.5cm, right=2.5cm, top=2.5cm, bottom=2.5cm}
\definecolor{fach}{HTML}{b35c00}
\definecolor{gray}{HTML}{666666}
\definecolor{successgreen}{HTML}{2a7a4b}
\definecolor{warnred}{HTML}{c62828}

\tcbuselibrary{skins,breakable}
\newtcolorbox{scorebox}{ colback=white, colframe=fach, fonttitle=\bfseries, title=Didaktik-Score, breakable }

\pagestyle{fancy}
\fancyhf{}
\fancyhead[L]{\small\textcolor{gray}{Digitale Grundlagen -- 009 Datenschutz}}
\fancyhead[R]{\small\textcolor{gray}{LE Tier 2 / Profil S}}
\fancyfoot[C]{\small\thepage{} / \pageref{LastPage}}
\renewcommand{\headrulewidth}{0.4pt}

\begin{document}

\begin{titlepage}
    \centering
    \vspace*{2cm}
    {\Large\textcolor{gray}{Digitale Grundlagen -- Senioren 65+}}\\[2cm]
    {\Huge\textcolor{fach}{\textbf{Langentwurf}}}\\[0.5cm]
    {\large Tier 2 | Profil S}\\[2cm]
    {\LARGE\textbf{009 -- Datenschutz \& Privatsphäre}}\\[0.5cm]
    {\large Modul B: Verbindungen \& Sicherheit}\\[2cm]
    \begin{tabular}{rl}
        \textbf{Zeitrahmen:} & 60--75 Minuten \\
        \textbf{Vorkenntnisse:} & iPhone-Einstellungen, Apps verstehen \\
    \end{tabular}
    \vfill
    {\normalsize Stand: \today}
\end{titlepage}

\tableofcontents
\newpage

\section{Bedingungsanalyse}

\subsection{Ausgangssituation}
\begin{itemize}
    \item Datenschutz ist abstrakt und schwer greifbar
    \item Viele akzeptieren alles ohne zu lesen (``Stimme zu'')
    \item Unklarheit: Was sammeln Apps/Firmen wirklich?
    \item Angst vor Überwachung, aber Unsicherheit was zu tun ist
\end{itemize}

\subsection{Typische Fragen}
\begin{itemize}
    \item ``Hört mein Handy mit?''
    \item ``Warum will die App meinen Standort wissen?''
    \item ``Was passiert mit meinen Daten?''
    \item ``Kann ich das abstellen?''
\end{itemize}

\section{Sachanalyse}

\subsection{Kernkonzepte}
\begin{enumerate}
    \item \textbf{Daten als Währung:} Kostenlose Apps verdienen Geld mit deinen Daten
    \item \textbf{Berechtigungen:} Apps fragen um Erlaubnis (Kamera, Mikrofon, Standort)
    \item \textbf{Einstellungen kontrollieren:} Du kannst Berechtigungen jederzeit ändern
    \item \textbf{Datensparsamkeit:} Nur das teilen, was wirklich nötig ist
    \item \textbf{Apple vs. Andere:} Apple betont Datenschutz (Tracking-Transparenz)
\end{enumerate}

\subsection{Alltagsanalogien}
\begin{tabular}{|l|p{8cm}|}
\hline
\textbf{Konzept} & \textbf{Analogie} \\
\hline
Datensammlung & Jemand folgt dir im Supermarkt und schreibt auf, was du kaufst \\
Berechtigungen & Schlüssel zu verschiedenen Räumen deines Hauses \\
Standort-Tracking & GPS-Sender an deinem Auto \\
Cookies & Fußspuren, die du im Internet hinterlässt \\
Datensparsamkeit & Nur die Schlüssel ausgeben, die wirklich gebraucht werden \\
\hline
\end{tabular}

\section{Didaktische Analyse}

\subsection{Stellung in der Reihe}
Konzeptueller Übergang: Von konkreter Technik (007, 008) zu Bewusstsein und Verhalten.

\subsection{Didaktische Reduktion}
\textbf{Ausgeklammert:} DSGVO-Details, Cookie-Banner im Detail, VPN, Verschlüsselung.

\textbf{Fokus:} Berechtigungen auf dem iPhone prüfen und verwalten.

\section{Lernziele}

\begin{tabular}{|c|p{7cm}|c|c|}
\hline
\textbf{Nr.} & \textbf{Die Teilnehmer können...} & \textbf{Stufe} & \textbf{Niveau} \\
\hline
FZ1 & erklären, warum Apps Daten sammeln wollen & Verstehen & Basis \\
FZ2 & die App-Berechtigungen auf ihrem iPhone finden & Anwenden & Basis \\
FZ3 & einer App den Standortzugriff entziehen & Anwenden & Basis \\
FZ4 & das App-Tracking-Fenster verstehen und entscheiden & Verstehen & Basis \\
FZ5 & den Datenschutzbericht lesen & Anwenden & Vertiefung \\
\hline
\end{tabular}

\section{Verlaufsplanung}

\begin{longtable}{|p{1cm}|p{2cm}|p{5cm}|p{3cm}|p{2cm}|}
\hline
\textbf{Zeit} & \textbf{Phase} & \textbf{Inhalt} & \textbf{TN-Aktivität} & \textbf{Material} \\
\hline
0--10 & Einstieg & ``Warum ist Facebook kostenlos?'' -- Daten als Währung & Nachdenken & -- \\
\hline
10--25 & Konzept & Was wird gesammelt? Warum? Wer verdient daran? & Zuhören, Fragen & S.1 \\
\hline
25--40 & Berechtigungen & Einstellungen $\rightarrow$ Datenschutz, Standort, Kamera prüfen & Am Gerät schauen & S.2 \\
\hline
40--55 & Tracking & App-Tracking-Transparenz erklären, ``Nicht erlauben'' zeigen & Mitmachen & S.2 \\
\hline
55--70 & Übungen & Berechtigungen einer App ändern, Checkliste & Bearbeiten & S.3 \\
\hline
\end{longtable}

\section{Lösungen}

\textbf{Wo finde ich die Berechtigungen?}
\begin{itemize}
    \item Einstellungen $\rightarrow$ Datenschutz \& Sicherheit $\rightarrow$ [Bereich] (z.B. Standortdienste)
    \item Oder: Einstellungen $\rightarrow$ [App-Name] $\rightarrow$ Berechtigungen
\end{itemize}

\textbf{Was bedeutet ``App-Tracking''?}
\begin{itemize}
    \item Apps wollen dein Verhalten über verschiedene Apps hinweg verfolgen
    \item ``Erlauben'' = App darf Werbeprofil erstellen
    \item ``Nicht erlauben'' = Empfohlen für mehr Privatsphäre
\end{itemize}

\section{Didaktik-Score}

\begin{scorebox}
\begin{tabular}{|c|l|c|c|p{3.5cm}|}
\hline
\# & Dimension & Gew. & Score & Notiz \\
\hline
1 & Differenzierung & 15\% & 9/10 & Basis + Datenschutzbericht als Vertiefung \\
2 & Taxonomie & 10\% & 9/10 & Verstehen + Anwenden \\
3 & Alltagsrelevanz & 10\% & 10/10 & Jeder ist betroffen \\
4 & Aktivierung & 10\% & 9/10 & Eigene Einstellungen prüfen \\
5 & Scaffolding & 10\% & 9/10 & Gute Analogien \\
6 & Feedback & 10\% & 9/10 & Ermutigung: ``Du hast die Kontrolle'' \\
7 & Sprachliche Klarheit & 10\% & 9/10 & Komplexes Thema vereinfacht \\
8 & Motivation & 10\% & 9/10 & Kontrolle zurückgewinnen \\
9 & Barrierefreiheit & 10\% & 9/10 & Klare Pfade \\
10 & Praktikabilität & 5\% & 9/10 & 70 Min angemessen \\
\hline
\multicolumn{3}{|r|}{\textbf{GESAMT:}} & \textbf{9.05/10} & \textcolor{successgreen}{\textbf{FREIGABE}} \\
\hline
\end{tabular}
\end{scorebox}

\end{document}
